\documentclass[10pt,a4paper]{amsart}
\usepackage[margin=19mm]{geometry}
\usepackage{amsmath,amssymb,mathtools}
\usepackage[scr=rsfs,cal=euler]{mathalfa}
\usepackage{xfrac}
\usepackage[unicode,hidelinks,bookmarks=false]{hyperref}
\newcommand{\ie}{i.e.\ }
\newcommand{\cf}{cf.\ }
\newcommand{\resp}{respectively\ }
\newcommand{\Subsection}{\subsection}
\providecommand{\nobreakdash}{\nobreak\hbox{-}}
\setlength{\emergencystretch}{2em}
\multlinegap0pt
\usepackage{xcolor}
\usepackage{tikz}
\newtheorem{lemma}{Lemma}
\newcommand{\eps}{\varepsilon}
\title{\textbf{Algorithmic List Decoding of Reed--Solomon Codes\\ up to Capacity}}
\author{Joshua Brakensiek, Yeyuan Chen, Aaron Putterman, Zihan Zhang, Kai Zhe Zheng}
\date{Source: arXiv:2609.08005v1 (CC BY 4.0)}
\begin{document}
\maketitle

\noindent\emph{Source and licence.} \textbf{Algorithmic List Decoding of Reed--Solomon Codes\\ up to Capacity}. Version arXiv:2609.08005v1 (CC BY 4.0), distributed by arXiv under the Creative Commons Attribution 4.0 International licence, http://creativecommons.org/licenses/by/4.0/, as recorded in the arXiv licence metadata for this exact version and shown on the abstract page https://arxiv.org/abs/2609.08005v1. Full licence notice: \texttt{licenses/arxiv-2609.08005-CC-BY-4.0.txt}.\par\smallskip

\noindent\textit{Editorial scope.} Counted unit: the article's lemma lem:global-budgets (source0.tex lines 542--567). The article's complete original proof of it is delivered with it (source0.tex lines 568--694, one \texttt{proof} environment of 127 lines). No mathematics is written, repaired or re-derived: the statement and the proof are the article's own lines, in the article's own notation, and no retained line is substituted or re-flowed.  Retained uncounted as the source-local context the proof needs: lines 357--364.  Nothing was omitted from any retained span, so the package carries no omission record.  No retained line contains a citation, so the wrapper carries no bibliography and no bibliography entry.  No retained line contains a figure environment, an image-inclusion command or drawing code, so no image is delivered and the package declares no asset dependency.  Editorial formatting only, in addition: an AMS \texttt{amsart} preamble that restates the article's own declarations and macros used by the retained lines, namely the environments \texttt{lemma} on separate counters, together with the standard packages those lines need.  The retained environments print in this document as Lemma 1 in order of appearance, whereas the article prints the same items with the numbering of its own full document.  The complete native source of the article as distributed in the arXiv e-print for this exact version is bundled unchanged as \texttt{sources/209685/source0.tex}.
\begin{equation} \label{eq:global_params}
A = \left\lceil
\eps n
\right\rceil, \qquad
d=\left\lceil \eps^{-3/\theta} \right \rceil,
\qquad
m=d^3,
\end{equation}

\begin{lemma}\label{lem:global-budgets}
The interpolation space $\mathcal Q$ satisfies the following properties:
\begin{itemize}
    \item Every $Q\in\mathcal Q$ satisfies
    \[
    \deg_{Y_i}Q
    \le B,
    \]
    for each $i \in \{0,\ldots, d\}$.
    \item For every $Q\in\mathcal Q$ and every
    $P\in\mathbb F_q[X]$ with $\deg P\le k-1$,
    \[
    \deg_X
    Q\bigl(X,P(X),P^{[1]}(X),\ldots,P^{[d]}(X)\bigr)
    <mA.
    \]

    \item The dimension of $\mathcal Q$ satisfies
    \[
    \dim\mathcal Q
    >
    \frac{\theta^3}{384}
    |\mathcal B|(k-1)m^3.
    \]
\end{itemize}
\end{lemma}

\begin{proof}
The first item follows directly from the definition of $\mathcal Q$.
Indeed, every monomial
\[
X^aY_0^{b_0}Y_1^{b_1}Y^c
\]
in $\mathcal Q$ satisfies
\[
b_0+b_1+|c|\le B.
\]
Hence the exponent of each variable $Y_i$, for $i\in\{0,\ldots,d\}$, is at
most $B$, and therefore
\[
\deg_{Y_i}Q\le B
\qquad
\]
for every $i \in \{0,\ldots, d\}$.

The second item follows straightforwardly too. Fix $P\in\mathbb F_q[X]$ of degree at most $k-1$. Since $\deg P^{[j]}\le k-1$
for every $j$, one can check that
\[
X^a
P(X)^{b_0}
\bigl(P^{[1]}(X)\bigr)^{b_1}
\prod_{j=2}^{d}
\bigl(P^{[j]}(X)\bigr)^{c_j},
\]
has $X$-degree at most
\[
a+(k-1)(b_0+b_1+|c|)<mA.
\]
Hence, any $Q \in \mathcal{Q}$ has $\deg_X
Q\bigl(X,P,P^{[1]},\ldots,P^{[d]}\bigr)
<mA$.

It remains to lower bound $\dim\mathcal Q$. First, by the definition of $A$ in \eqref{eq:global_params} and the fact that $k \leq (1-\theta) \eps$, we have
\[
\frac{A}{k-1}
\ge
\frac{\epsilon n}{(1-\theta)\epsilon n}
=
\frac{1}{1-\theta}
\ge
1+\theta.
\]
On the other hand, for every $c\in\mathcal B$, the definition of
$\mathcal B$ gives
\[
|c|
\le
\left\lfloor
\left(1+\frac{3\theta}{4}\right)m
\right\rfloor
\le
\left(1+\frac{3\theta}{4}\right)m.
\]
Consequently,
\begin{equation}\label{eq:rank-lb-tmp}
m\frac{A}{k-1}-|c|
\ge
\frac{\theta m}{4}.
\end{equation}

Now fix $c\in\mathcal B$ and an integer
\[
0\le b_1\le
\left\lfloor\frac{\theta m}{4}\right\rfloor.
\]
For every pair $(a,b_0)\in\mathbb Z_{\ge0}^2$ satisfying
\[
a+(k-1)b_0
<
mA-(k-1)(|c|+b_1),
\]
observe that the corresponding monomial $X^aY_0^{b_0}Y_1^{b_1}Y^c$
belongs to $\mathcal Q$. Indeed, the weighted-degree condition holds by
construction, and since $a\ge0$,
\[
b_0+b_1+|c|
<
\frac{mA}{k-1}
\le
\left\lceil\frac{mA}{k-1}\right\rceil
=
B.
\]

The number of such pairs $(a,b_0)$ is at least
\[
\frac{
\bigl(mA-(k-1)(|c|+b_1)\bigr)^2
}{
2(k-1)
}.
\]
Using \eqref{eq:rank-lb-tmp}, this is at least
\[
\frac{k-1}{2}
\left(
\frac{\theta m}{4}-b_1
\right)^2.
\]
Therefore,
\[
\begin{aligned}
\dim\mathcal Q
&\ge
\frac{k-1}{2}
|\mathcal B|
\sum_{b_1=0}^{\lfloor\theta m/4\rfloor}
\left(
\frac{\theta m}{4}-b_1
\right)^2\\
&\ge
\frac{k-1}{2}
|\mathcal B|
\int_0^{\theta m/4}
\left(
\frac{\theta m}{4}-x
\right)^2\,dx\\
&=
\frac{\theta^3}{384}
|\mathcal B|(k-1)m^3.
\end{aligned}
\]
This proves the final item.
\end{proof}

\end{document}
